\section*{Hoofdstuk \ref{ch:g11:ocean-ridges} --- Oceanische ruggen en de oceanische lithosfeer}

\begin{solution}{exo:g11:ocean-ridges:1}
\omterm{def:g11:ocean-ridges:peridotite}{Peridotiet} is een donker, grofkorrelig gesteente dat vooral uit olivijn bestaat, met
wat pyroxeen; een groot deel van de mantel bestaat eruit. \omterm{def:g11:ocean-ridges:gabbro}{Gabbro} is een donker,
grofkorrelig stollingsgesteente van plagioklaas-veldspaat en pyroxeen, het meest
voorkomende gesteente van de diepe \omterm{def:g11:ocean-ridges:oceanic-lithosphere}{oceanische korst}. \omterm{def:g11:ocean-ridges:gabbro}{Gabbro} hoort bij de korst,
\omterm{def:g11:ocean-ridges:peridotite}{peridotiet} bij de mantel. \omterm{def:g11:ocean-ridges:gabbro}{Gabbro} kristalliseerde langzaam, diep onder de zeebodem,
waar het magma langzaam afkoelde: zijn kristallen hadden tijd om te groeien.
Basaltlava koelde snel af op de zeebodem, en haar kristallen bleven klein.
\end{solution}

\begin{solution}{exo:g11:ocean-ridges:2}
Van de zeebodem naar beneden: sedimenten (geen op de as), basaltlava's (kussens
en plaatvormige stromen), het gangencomplex, \omterm{def:g11:ocean-ridges:gabbro}{gabbro}, en dan, onder de \omterm{def:g11:seismic-structure:moho}{Moho}, de
\omterm{def:g11:ocean-ridges:peridotite}{peridotiet} van de mantel. Het gat doorkruiste $1507/7100 \approx 0.21$, ongeveer
een vijfde van de korst: de diepe \omterm{def:g11:ocean-ridges:gabbro}{gabbro}'s en de \omterm{def:g11:seismic-structure:moho}{Moho} zijn nog nooit door een boor
bereikt.
\end{solution}

\begin{solution}{exo:g11:ocean-ridges:3}
$d(0) = \qty{2500}{m}$; $d(9) = 2500 + 350 \times 3 = \qty{3550}{m}$;
$d(25) = 2500 + 350 \times 5 = \qty{4250}{m}$.
\end{solution}

\begin{solution}{exo:g11:ocean-ridges:4}
Een \omterm{def:g11:ocean-ridges:normal-fault}{afschuiving} is een hellende breuk waarlangs het blok boven het breukvlak
naar beneden is bewogen ten opzichte van het blok eronder. Bij een langzame rug
trekken de platen de koude, brosse \omterm{def:g11:lithosphere-asthenosphere:lithosphere}{lithosfeer} uit elkaar: ze wordt uitgerekt, en
ze breekt. De blokken boven de breuken glijden omlaag, naar de as, en het axiale
dal is het blok dat ertussen is weggezakt.
\end{solution}

\begin{solution}{exo:g11:ocean-ridges:5}
Het kookt niet door de druk: meer dan twee kilometer water drukt erop,
en onder die druk blijft water vloeibaar tot ver boven
\qty{100}{\celsius}. De zwarte ``rook'' is een wolk van piepkleine korrels
ijzersulfiden en andere mineralen, die neerslaan wanneer het hete bronwater, rijk
aan opgeloste metalen, het koude zeewater raakt.
\end{solution}

\begin{solution}{exo:g11:ocean-ridges:6}
(1) Op ongeveer \qty{54}{km} diepte, bij ongeveer \qty{1324}{\celsius}. (2) Op
ongeveer \qty{85}{km} diepte, bij ongeveer $1400 + 0.45 \times 85 \approx
\qty{1438}{\celsius}$. (3) De hetere mantel begint dieper te smelten, dus een
langere kolom opstijgend gesteente smelt voordat ze de korst bereikt: ze maakt meer
basalt, en een dikkere korst.
\end{solution}

\begin{solution}{exo:g11:ocean-ridges:7}
(1) $1489 - 1345 = \qty{144}{\celsius}$ onder zijn smelttemperatuur.
(2) Het stijgt $100 - 54 = \qty{46}{km}$; het verschil krimpt met
$144/46 \approx \qty{3.1}{\celsius}$ per kilometer gemiddeld. (3) Nee: onderweg
naar boven koelt het gesteente zelfs een beetje af, met ongeveer $0.45 \times 46 \approx
\qty{21}{\celsius}$. Het is zijn smelttemperatuur die daalt, veel sneller,
terwijl de druk afneemt.
\end{solution}

\begin{solution}{exo:g11:ocean-ridges:8}
De halve snelheden zijn \qty{7.3}{cm} en \qty{1.65}{cm} per jaar, d.w.z.\
\qty{73}{km} en \qty{16.5}{km} per miljoen jaar. Om aan één kant \qty{1000}{km}
bodem te maken: $1000/73 \approx 13.7$ miljoen jaar voor de Oost-Pacifische
Rug, $1000/16.5 \approx 61$ miljoen jaar voor de Midden-Atlantische Rug.
\end{solution}

\begin{solution}{exo:g11:ocean-ridges:9}
(1) \qty{5}{cm} is \qty{5e-5}{km}: $65000 \times 5 \times 10^{-5} =
\qty{3.25}{km^2}$ nieuwe bodem per jaar. (2) $3.25 \times 7.1 \approx
\qty{23}{km^3}$ korst per jaar. (3) $3.6 \times 10^8 / 3.25 \approx 1.1
\times 10^8$ jaar, ongeveer 110 miljoen jaar: in dit tempo zouden de ruggen de hele
oceaanbodem in ongeveer honderd miljoen jaar kunnen maken, een korte tijd
vergeleken met de ouderdom van de Aarde. Daarom is geen enkele oceaanbodem heel oud.
\end{solution}

\begin{solution}{exo:g11:ocean-ridges:10}
(1) Het breukvlak maakt \ang{60} met de horizontaal, en de daling is de
verticale zijde van een rechthoekige driehoek waarvan de schuine zijde de verschuiving is:
$1.5/\sin\ang{60} \approx \qty{1.73}{km}$. (2) Elke breuk beweegt de kanten
horizontaal uit elkaar met $1.5/\tan\ang{60} \approx \qty{0.87}{km}$; de twee
breuken samen met ongeveer \qty{1.73}{km}. (3) Nee. De \qty{25}{km} is de
afstand tussen de breuken die het weggezakte blok begrenzen, niet de hoeveelheid
uitrekking, die maar ongeveer \qty{1.7}{km} is; en de rug gaat ook verder open
door nieuwe gangen en lava's bij de as, niet alleen door breukvorming.
\end{solution}

\begin{solution}{exo:g11:ocean-ridges:11}
$d(25) = 6400 - 3200\,e^{-25/62.8} \approx 6400 - 3200 \times 0.672 \approx
\qty{4251}{m}$; $d(100) = 6400 - 3200\,e^{-100/62.8} \approx 6400 - 3200
\times 0.203 \approx \qty{5749}{m}$. Tussen 25 en 100 miljoen jaar zakt de
bodem ongeveer \qty{1500}{m} weg. De wortelwet geeft
$2500 + 350 \times 5 = \qty{4250}{m}$ bij 25 miljoen jaar: de twee wetten
komen overeen waar ze allebei gelden.
\end{solution}

\begin{solution}{exo:g11:ocean-ridges:12}
(1) $\alpha\,\Delta T\,L = 3.2 \times 10^{-5} \times 700 \times 100 =
\qty{2.24}{km}$. (2) $s = 3380 \times 2.24/(3380 - 1025) = 7571/2355 \approx
\qty{3.2}{km}$. (3) De exponentiële wet laat in totaal $6400 - 2500 =
\qty{3900}{m}$ wegzakken toe: afkoeling en krimp alleen verklaren ongeveer
vier vijfde ervan, met eenvoudige waarden voor de temperatuurdaling en de
dikte.
\end{solution}

\begin{solution}{exo:g11:ocean-ridges:13}
(1) $54 - 7 = \qty{47}{km}$ tegenover $85 - 7 = \qty{78}{km}$: de hetere
kolom is $78/47 \approx 1.66$ keer langer. (2) $7.1 \times 1.66 \approx
\qty{11.8}{km}$. (3) De gemeten \qty{10.3}{km} is dikker dan normaal, zoals
voorspeld, maar minder dan de schatting: de diepste kilometers van de kolom,
net boven het begin van het smelten, geven minder smelt dan de ondiepere, dus
de aanname van evenveel smelt per kilometer overschat.
\end{solution}

\begin{solution}{exo:g11:ocean-ridges:14}
(1) De wet geeft \qty{3178}{m}, \qty{4160}{m}, \qty{4782}{m} en
\qty{5267}{m}; de gemeten medianen zijn ondieper met \qty{64}{m},
\qty{142}{m}, \qty{317}{m} en \qty{619}{m}. (2) Het verschil groeit met de
ouderdom. (3) De wet geeft de diepte van de bodem zelf, gecorrigeerd voor zijn
sedimenten; de gemeten cellen dragen hun sedimenten, die dikker zijn op
oudere bodem, en bevatten onderzeese bergen en vulkanische plateaus, die
ondieper zijn dan de bodem eromheen.
\end{solution}

\begin{solution}{exo:g11:ocean-ridges:15}
(1) $p = 1025 \times 9.8 \times 2500 \approx \qty{25.1}{MPa}$, plus
\qty{0.1}{MPa}: ongeveer \qty{25.2}{MPa}. (2) Dat is boven de kritieke
druk, \qty{22.064}{MPa}: het bronwater kan niet koken, hoe heet het ook is.
(3) De druk bereikt \qty{22.064}{MPa} op de diepte
$(22.064 - 0.1) \times 10^6/(1025 \times 9.8) \approx \qty{2190}{m}$: elke
bron dieper dan ongeveer \qty{2200}{m} kan niet koken.
\end{solution}

\begin{solution}{pb:g11:ocean-ridges:1}
\textbf{1.} $14.6/2 = \qty{7.3}{cm}$ per jaar, d.w.z.\ \qty{73}{km} per miljoen
jaar.

\textbf{2.} Op de diepte van de as van de rug, $d(0) = \qty{2500}{m}$.

\textbf{3.} Van de zeebodem naar beneden: basaltlava's (kussens en plaatvormige stromen),
het gangencomplex en \omterm{def:g11:ocean-ridges:gabbro}{gabbro}, in totaal ongeveer \qty{7.1}{km} korst; onder de
\omterm{def:g11:seismic-structure:moho}{Moho} de \omterm{def:g11:ocean-ridges:peridotite}{peridotiet} van de mantel, tot de onderkant van de \omterm{def:g11:lithosphere-asthenosphere:lithosphere}{lithosfeer}, die bij de as
nog heel dun is. Later zullen er sedimenten op neerslaan.

\textbf{4.} Wanneer de platen uit elkaar bewegen, stijgt de \omterm{def:g11:lithosphere-asthenosphere:lithosphere}{asthenosfeer} onder de
rug op. Haar temperatuur daalt maar met ongeveer \qty{0.45}{\celsius} per
kilometer, terwijl haar smelttemperatuur veel sneller daalt naarmate de druk
afneemt. Op ongeveer \qty{54}{km} diepte ontmoeten de twee elkaar: het gesteente begint
te smelten, en smelt verder terwijl het opstijgt, zonder dat er warmte wordt
toegevoegd (\omterm{def:g11:ocean-ridges:partial-melting}{decompressiesmelten}).

\textbf{5.} Elk miljoen jaar maakt de rug $2 \times 73 = \qty{146}{km}$
bodem aan zijn twee kanten; per kilometer rug is het volume korst
$146 \times 7.1 \times 1 \approx \qty{1040}{km^3}$.

\textbf{6.} $73 \times t$: \qty{73}{km}, \qty{292}{km}, \qty{1168}{km} en
\qty{4672}{km}.

\textbf{7.} $2000/73 \approx 27$ miljoen jaar.

\textbf{8.} Platen veranderen over tientallen miljoenen jaren van snelheid en
richting, zoals de knik van de Hawaii-Keizerketen laat zien
(\cref{ch:g11:plate-motions}); de snelheid van vandaag hoeft niet 64
miljoen jaar te gelden.

\textbf{9.} Verder van de as ligt de korst niet meer boven een magmalens: ze
koelt af, en er is eronder niet meer genoeg warmte om de circulatie aan te drijven
en het water te verhitten.

\textbf{10.} Ten westen van de Oost-Pacifische Rug reist het stuk op de Pacifische
plaat. Naar het westen eindigt de Pacifische plaat langs diepe troggen, waar ze
onder haar buren wegzakt: subductiezones (\cref{ch:g11:subduction-zones}).

\textbf{11.} $d(1) = \qty{2850}{m}$, $d(4) = 2500 + 700 = \qty{3200}{m}$,
$d(16) = 2500 + 1400 = \qty{3900}{m}$, $d(64) = 2500 + 2800 = \qty{5300}{m}$.

\textbf{12.} $3900 - 2500 = \qty{1400}{m}$ in de eerste 16 miljoen jaar,
en $5300 - 3900 = \qty{1400}{m}$ tussen 16 en 64 miljoen jaar: hetzelfde
wegzakken, in een drie keer zo lange tijd. De bodem zakt steeds langzamer weg
naarmate hij ouder wordt.

\textbf{13.} $350\sqrt{t} = 1500$ geeft $\sqrt{t} \approx 4.29$, dus
$t \approx 18.4$ miljoen jaar.

\textbf{14.} $6400 - 3200\,e^{-64/62.8} \approx 6400 - 3200 \times 0.361
\approx \qty{5245}{m}$, ongeveer \qty{55}{m} minder dan de wortelwet: de
twee wetten komen tot op ongeveer \qty{1}{\percent} overeen.

\textbf{15.} De gemeten mediaan is ongeveer \qty{650}{m} ondieper dan de
wet. De wet beschrijft de bodem gecorrigeerd voor het gewicht van zijn sedimenten;
de gemeten diepten bevatten de sedimenten, vele tientallen miljoenen jaren
ervan op zulke bodem, en de onderzeese bergen en vulkanische plateaus, die allemaal
de gemeten bodem ondieper maken.

\textbf{16.} $\alpha\,\Delta T = 3.2 \times 10^{-5} \times 700 = 0.0224$: het
manteldeel is \qty{2.24}{\percent} dichter,
$3380 \times 1.0224 \approx \qty{3456}{kg/m^3}$.

\textbf{17.} Het manteldeel is $100 - 7.1 = \qty{92.9}{km}$ dik. Gemiddelde
dichtheid: $(7.1 \times 2900 + 92.9 \times 3456)/100 \approx (20590 + 321060)/100
\approx \qty{3416}{kg/m^3}$.

\textbf{18.} $3416 > 3380$: de oude \omterm{def:g11:lithosphere-asthenosphere:lithosphere}{lithosfeer} als geheel is ongeveer
\qty{36}{kg/m^3}, ongeveer \qty{1}{\percent}, dichter dan de \omterm{def:g11:lithosphere-asthenosphere:lithosphere}{asthenosfeer}
eronder. De lichte korst is niet meer genoeg om haar koude,
dichte manteldeel te compenseren.

\textbf{19.} Omdat ze dichter is dan de \omterm{def:g11:lithosphere-asthenosphere:lithosphere}{asthenosfeer}, kan ze bij een trog terug in de
mantel zakken: subductie, het onderwerp van het volgende hoofdstuk.

\textbf{20.} Bij 64 miljoen jaar ligt het stuk bodem op ongeveer
\qty{5300}{m}: ongeveer \textbf{5.3~km} diep, \qty{2.8}{km} dieper dan waar het
geboren werd.
\end{solution}
